% Generated draft: requires mathematical and source review.
\documentclass[11pt]{article}
\usepackage[margin=25mm]{geometry}
\usepackage{fontspec}
\setmainfont{Latin Modern Roman}
\usepackage{amsmath,amssymb,mathtools,unicode-math}
\setmathfont{Latin Modern Math}
\usepackage{enumitem,needspace,xurl,hyperref}
\setlength{\parindent}{0pt}
\setlength{\parskip}{5pt}
\setlength{\emergencystretch}{3em}
\newcommand{\N}{\mathbb{N}}
\newcommand{\Z}{\mathbb{Z}}
\newcommand{\Q}{\mathbb{Q}}
\newcommand{\R}{\mathbb{R}}
\newcommand{\C}{\mathbb{C}}
\newcommand{\F}{\mathbb{F}}
\newcommand{\A}{\mathbb{A}}
\newcommand{\PP}{\mathbb{P}}
\newcommand{\E}{\mathbb{E}}
\newcommand{\eps}{\varepsilon}
\newcommand{\dd}{\,\mathrm{d}}
\begin{document}
\section*{UnsolvedMath ID: 30000654}
\textbf{Source problem code:} catalogue entry
\subsection{Definitions and mathematical statement}
Let $p,q\in\Z[x,y]$ with $p(0,0)=q(0,0)=0$, and let $f=(p,q)$ be the induced polynomial self-map of the affine plane. Dominance is understood on the characteristic-zero generic fiber: the image of $f:\mathbb A^2_{\Q}\to\mathbb A^2_{\Q}$ is Zariski dense. Write $f^m$ for the $m$-fold iterate. A point $a$ is periodic if $f^m(a)=a$ for some integer $m\geq1$, including fixed points.

For each prime $\ell$, let $\overline f_\ell:\F_\ell^2\to\F_\ell^2$ be coefficientwise reduction. Preserve the exact finite-field hypothesis
\[
\forall\ell\text{ prime}\ \exists a_\ell\in\F_\ell^2\setminus\{(0,0)\}\ \exists m_\ell\geq1,
\qquad \overline f_\ell^{\,m_\ell}(a_\ell)=a_\ell.
\]
The prime-dependent points and periods need not be compatible or uniformly bounded, and the hypothesis says every prime, not only almost all primes.

The literal requested conclusion is
\[
\exists a\in\Q^2\ \exists m\geq1:\quad f^m(a)=a.
\]
However, the imposed equations already give $f(0,0)=(0,0)$, so the rational origin satisfies this conclusion without any finite-field or dominance hypothesis. The nonzero-point exclusion occurs only in the finite-field antecedent in both supplied formulations. A potentially substantive variant would require $a\ne(0,0)$ also over $\Q$, but that stronger conclusion is not written in the source and is not silently substituted. Clarification of this missing exclusion, or another intended nontriviality requirement, is essential to recover an open problem.
\subsection{Short English statement}
The source asks whether nonzero periodic points modulo every prime force a rational periodic point for a dominant plane map fixing the origin. As written, the rational origin is already periodic; a nonzero rational-point requirement would be a stronger, unstated question.
\subsection{Sources}
\begin{itemize}
\item Supplied report: \emph{Affine Algebraic Geometry} (2007), Oberwolfach Reports, DOI 10.4171/owr/2007/01, pp. 5--82; proposed year 2007. \url{https://doi.org/10.4171/owr/2007/01}.
\item Local UnsolvedMath registry, permanent ID 30000654, catalogue entry; created 25 August 2026, updated 16 September 2026. Its corrected\_verified note describes balancing mathematical notation. Both original and clean statements exclude the origin only over finite fields. The supplied 21 August 2026 literature triage labels the extracted claim open and reports no located authoritative complete resolution; this label does not resolve the literal fixed-origin issue. No additional external reference is supplied; no independent literature assessment is made here.
\end{itemize}
\end{document}
